%%=====================================================================
%% M2 — WP3 / R3.1–R3.6 (RE-SCOPED, 8 Nov 2026): the field equation on M₁₂.
%% SURGICAL FRAGMENT (insertion only; no preamble). UK English. "We".
%% Under PI rulings Q1 (one gravitational term), Q8 (Riemannian; time =
%% flow parameter), Q12 (induced Spin^c connection, no gauge coupling), and
%% the re-scoped R3.1 (Φ_QC a dilaton on M₁₂ with coupling ζ; the warp is
%% the geometry's response, a = −ζΦ; Paper 2's 2D system is the base
%% restriction). Supersedes the withdrawn 4D-radion draft entirely (SOW §8).
%% Depends on WP0 (def:scheme, def:riemannian, def:omega, prop:parity) and
%% WP2 (sec:handedness). Placeholder: Section~\ref{sec:warp} = WP4 (warp exponent).
%%=====================================================================

\section{The field equation}\label{sec:equation}

The purpose of this section is to write the action on $M_{12}$ that the
axioms admit, to vary it, and to display the field equation that results.
In plain terms: we write down the one number that every configuration of
the shape and its fields is assigned, allow it to contain only what the
axioms force or the smallest option they leave open, and ask for the
configurations at which it does not change. The answer is Einstein's
equation with one additional term, and the section says which axiom is
responsible for each part: closure for the curvature and the vacuum term,
order for the matter and the force field, growth and decay for the
additional term. The equation is twelve-dimensional. Its restriction to
the base of $\Keight$ is the system in which Section~\ref{sec:warp} derives
the warp exponent, and its four-dimensional consequences are the subject
of later papers of the series; neither is needed to state it. We announce
the result.

\begin{theorem}[the Master Equation]\label{thm:master}
Under the action \eqref{eq:action} the metric of $M_{12}$ satisfies
\begin{equation}\label{eq:master}
\begin{gathered}
  G_{MN}\;=\;8\pi G_{\mathrm{eff}}\,T_{MN}\;+\;\zeta\,T^{\QCD}_{MN}\;+\;\Lambda_{\mathrm{eff}}\,g_{MN}\,,\\
  T^{\QCD}_{MN}\;\equiv\;\langle\nabla_{M}\QCD,\nabla_{N}\QCD\rangle_{M}-\tfrac12\,g_{MN}\,\langle\nabla\QCD,\nabla\QCD\rangle_{M}\,,
\end{gathered}
\end{equation}
where $\QCD=\omega\lvert\psi-\bar\psi\rvert^{2}$ is the dilaton of
Definition~\ref{def:omega}, a section of the Clifford module $M\cong\Hquat^{32}$,
$\langle u,v\rangle_{M}=\sum_{\alpha}\mathrm{Re}(u_{\alpha}\bar v_{\alpha})$ the
module inner product, $\zeta$ the coupling, $T_{MN}$ the conserved
stress tensor of the spinor sector of A3 together with the field strength
of the $\Spinc$ line, $8\pi G_{\mathrm{eff}}=\kappa_{12}^{2}$, and
$\Lambda_{\mathrm{eff}}$ collects the vacuum terms. The dilaton satisfies
\begin{equation}\label{eq:dilaton}
  \zeta\,\Box\QCD\;=\;V'(\QCD)\,,
\end{equation}
so that $\Box\QCD=0$ wherever the potential is stationary or absent.
\end{theorem}

In words, the tensor $T^{\QCD}_{MN}$ is built as follows. Take the rate of
change of the dilaton in direction $M$ and the rate of change in direction
$N$; these are each thirty-two quaternions, and their inner product over
the module, summing the real parts of the products component by component,
is a single real number $\langle\nabla_{M}\QCD,\nabla_{N}\QCD\rangle_{M}$. That is
the first term: the stress the dilaton's variation puts between directions
$M$ and $N$. The second term subtracts half the total squared rate of change
of the dilaton, in all directions at once, spread evenly over the metric
$g_{MN}$; it is the pressure the dilaton's gradient exerts isotropically. The
difference is the standard stress tensor of a scalar field, familiar from
every textbook, with the one change that the squares are inner products
over the module. For a single real component it is exactly
$\nabla_{M}\QCD\nabla_{N}\QCD-\tfrac12g_{MN}\lvert\nabla\QCD\rvert^{2}$.

Read term by term, \eqref{eq:master} is Einstein's equation, $G_{MN}=8\pi G\,T_{MN}+\Lambda g_{MN}$,
with the one additional term $\zeta T^{\QCD}_{MN}$. The curvature on the
left and the vacuum term exist because the shape closes on itself (A1);
the matter and force term $T_{MN}$ exists because order matters (A3) and
the algebra of order acts on spinors; the additional term exists because
the shape grows and decays as a function of itself (A2), its field being
written on the two sheets of that law; and the two sheets are two, rather
than one sheet and its reflection, because the shape never stands still
(A4, Section~\ref{sec:handedness}). Einstein's equation is the case in
which the second and fourth axioms are absent.

Two readings are fixed now. By Definition~\ref{def:riemannian} the metric
is positive-definite and every ``time'' statement below is a statement
along the flow. And $\zeta$ in \eqref{eq:master} is a coupling constant of
the twelve-dimensional theory; Section~\ref{sec:warp} derives its value, $12\pi$, as the
exponent by which the warp of A2 falls across a half-period of the base
when the geometry responds to the dilaton's profile. The number is one
number seen in two places, and the equation displays it once.
%% Re-scoped R3.1. The withdrawn 4D radion computation is not this object.

\subsection{The action}\label{sec:action}

\begin{definition}[the action on $M_{12}$]\label{def:action}
\begin{equation}\label{eq:action}
\begin{aligned}
  S\;=\;\frac{1}{2\kappa_{12}^{2}}\int_{M_{12}}\!\sqrt{g}\,
  \Bigl[\,R\;-\;2\Lambda_{12}\;-\;\tfrac14\,F_{MN}F^{MN}
  \;-\;&\zeta\,\langle\nabla\QCD,\nabla\QCD\rangle_{M}\;-\;2V(\QCD)\Bigr]\\
  &+\;S_{\mathrm{spinor}}[g,\Psi,A]\,,
\end{aligned}
\end{equation}
with $F=\mathrm{d}A$ the curvature of the induced connection on the
determinant line $L=K^{-1}$ (Section~\ref{sec:connection}), $\QCD$ the
dilaton of Definition~\ref{def:omega}, $V$ its potential, and
$S_{\mathrm{spinor}}$ the Dirac action of the $\Spinc$ spinor of A3 on the
bundle of M1 Proposition~7.3.
\end{definition}
%% CONVENTION RULED (PI, 9 Nov 2026, M2-7): the Φ kinetic term sits inside
%% the 1/2κ² bracket with coefficient ζ; this reproduces the canonical display
%% with no κ² on T^Φ. CFM v12 l.1525's −½ζ outside the bracket is the
%% superseded normalisation (Corpus Errors entry 20).

In words: the number $S$ assigned to a configuration is the integral over
the whole twelve-dimensional shape of five contributions, all weighted by
the volume element $\sqrt{g}$ and divided by twice the gravitational
constant $\kappa_{12}^{2}$. The first, $R$, is the scalar curvature --- how
much the shape bends at each point. The second, $-2\Lambda_{12}$, is the
energy of empty volume, a constant per unit volume. The third,
$-\tfrac14F_{MN}F^{MN}$, is the energy of the field strength of the
$\Spinc$ line, the one abelian field the geometry carries. The fourth,
$-\zeta\langle\nabla\QCD,\nabla\QCD\rangle_{M}$, is the energy stored in the
dilaton's variation from point to point, the module inner product of its
gradient with itself, multiplied by the coupling $\zeta$. The fifth,
$-2V(\QCD)$, is the potential energy of the dilaton's value. Added to all
of these is the action of the spinor field of the third axiom, the Dirac
action $S_{\mathrm{spinor}}$, which depends on the metric, the spinor
$\Psi$ and the connection $A$. Each term is accounted for by an axiom or a
reading, and we go through them.

\paragraph{One gravitational term.}
By A1 the domain is compact without boundary (M1 Definition~2.1), so the
action has no boundary terms and its stationary points are Euler--Lagrange
equations with no boundary data \cite{Wald}. The divergence-free symmetric
tensors built from the metric and two derivatives in twelve dimensions
form the five-term Lovelock family \cite{Lovelock1971}; every term beyond
the first carries an independent dimensionful coefficient, an unforced
datum under criterion~(1) of M1 Definition~3.2, excluded on the ground that
excluded the twisted $\Spinc$ lines in M1 Proposition~7.3. The Einstein
term survives because without it there is no equation; its coefficient is
the declared scale of M1 Theorem~8.1. Once one term linear in second
derivatives is admitted its uniqueness holds in every dimension
\cite{Vermeil1917,Weyl1922,Cartan1922}. \rigD
%% PI ruling Q1.

\paragraph{The connection of the line.}\label{sec:connection}
M1 Proposition~7.3 selects the $\Spinc$ structure induced by the
almost-complex structure of A3's selected unit, with determinant line
$L=K^{-1}$ (M1 Definition~2.5). The connection on $L$ is part of the
geometry's own data, and its field strength is the one abelian field
strength the geometry carries. Its energy therefore enters
\eqref{eq:action} inside the gravitational bracket, weighted by
$1/2\kappa_{12}^{2}$: a separate gauge coupling for a connection the
geometry already owns would be an unforced datum, excluded by
criterion~(1). What is \emph{not} fixed by the metric is the flux: the
quanta of $F$ through the cycles of $\Keight$ are topological data, with
total $N=\chi(\Keight)$ by the tadpole condition (Section~\ref{sec:warp}), and they sit
on the base cycle because it is the only cycle with cone points at which a
localised flux can satisfy the first-order equation of Section~\ref{sec:warp}. \rigD
%% PI rulings Q12/Q14 (weak (c), 8 Nov 2026). Consistent with CFM v12 l.1419
%% (Maxwell term inside the bracket) and Paper I Addendum §4.4–4.6 (κ²_{2,eff}
%% = κ₈²/(Vol₀ e^{6A₀}); twelve units on the base cycle). The identification
%% F = Ricci form, briefly adopted, is withdrawn (Minimality Register M2-2).

\paragraph{The dilaton term, and why the sheets enter through $\QCD$ alone.}
The fundamental fields are the sheets $\psi,\bar\psi$ and the Coherent
Field $\omega$ (Section~\ref{sec:variables}); the spinor's own dynamics are
the Dirac action in $S_{\mathrm{spinor}}$. The dilaton
$\QCD=\omega\lvert\psi-\bar\psi\rvert^{2}$ is composite and module-valued, so a
two-derivative term in it is a particular interaction among the
fundamental fields and their gradients, and writing it as the only such
term is a statement: the geometry couples to the sheets through $\QCD$ and
through nothing else. Which quadratic form in $\nabla\QCD$ appears is not a
further choice: the invariant symmetric forms on an irreducible
quaternionic-unitary module form a one-dimensional space, so covariance
under the fibre algebra forces the module inner product
$\langle\nabla_{M}\QCD,\nabla_{N}\QCD\rangle_{M}$ --- the trace of the square,
not the square of the trace. For a vacuum symmetric under the module the
trace is $32$ times a single component, and the three readings of
``$\zeta\nabla\QCD\nabla\QCD$'' for a module-valued field --- one component,
the trace, the square of the trace --- carry the weights $1$, $32$ and $1024$
(\texttt{L7\_three\_readings}); only the second is covariant. \rigR
We adopt this by the standing principle of Section~\ref{sec:readings}, and
the reason it is the minimal coupling is Proposition~\ref{prop:parity}. If curvature is to respond to the sheets
through a scalar that is even under the fold and vanishes both where the
sheets coincide and where growth stands against decay, $\QCD$ is the unique
lowest-order candidate, and $\lvert\nabla\QCD\rvert^{2}$ is then the unique
two-derivative term with those properties; the competing terms fail the
vanishing conditions ($\lvert\nabla\omega\rvert^{2}$ does not vanish at
$\psi=\bar\psi$; $\lvert\nabla\psi\rvert^{2}$ vanishes nowhere). This is
A2 read on the metric: the geometry grows and decays in response to the
asymmetry of the sheets weighted by their coherence, and to nothing else.
$\zeta$ is the strength of that one coupling and is never absorbed into a
redefinition of the field. \rigD\ (minimality)
%% The corpus posits S_Φ (CFM l.1525; Paper 2 §7.3 "microscopic origin a
%% separate question"). M2 names the step and gives the parity theorem as
%% its reason. Tagged [D]: minimal, not the only conceivable, coupling.

\paragraph{The potential and the vacuum term.}
$\Lambda_{12}$ is the energy of empty twelve-volume; $V(\QCD)$ is the
dilaton's potential, whose form on the base is obtained in Section~\ref{sec:warp} from
the fibre curvature and $\Lambda_{12}$ when the action is evaluated on the
warped configuration. Neither is a free parameter of this paper:
$\Lambda_{12}$ is fixed by the requirement that the base system of
Section~\ref{sec:warp} admit a solution with the given topology, and $V$ is derived
there. Earlier work in the programme took $\Lambda_{12}=0$ as an axiom of
the construction; in this series that statement is the case in which the
solvability condition returns zero, to be established or refuted where the
fibre radii are derived, and it is not assumed here. \rigD
%% PI ruling (Λ₁₂ by stationarity, 8 Nov 2026).

\subsection{Variation}\label{sec:variation}

\begin{proof}[Proof of Theorem~\ref{thm:master}]
Varying \eqref{eq:action} with respect to $g^{MN}$,
\begin{equation}\label{eq:var}
\begin{aligned}
  G_{MN}+\Lambda_{12}g_{MN}
  \;=\;&\zeta\Bigl(\langle\nabla_{M}\QCD,\nabla_{N}\QCD\rangle_{M}-\tfrac12 g_{MN}\langle\nabla\QCD,\nabla\QCD\rangle_{M}\Bigr)
  \;-\;V(\QCD)\,g_{MN}\\
  &+\;\tfrac12\Bigl(F_{MP}F_{N}{}^{P}-\tfrac14 g_{MN}F^{2}\Bigr)
  \;+\;\kappa_{12}^{2}\,T^{\mathrm{sp}}_{MN}\,,
\end{aligned}
\end{equation}
where $T^{\mathrm{sp}}_{MN}=-\frac{2}{\sqrt g}\frac{\delta S_{\mathrm{spinor}}}{\delta g^{MN}}$.
The first bracket is $T^{\QCD}_{MN}$ with coefficient $\zeta$ and nothing
else; it is the canonical stress tensor of a module-valued scalar, the
trace over $M$ of the scalar form, and its form is forced by the kinetic
term, not chosen. For a single component it reduces to
$\nabla_{M}\QCD\nabla_{N}\QCD-\tfrac12g_{MN}\lvert\nabla\QCD\rvert^{2}$, which is
the display the programme has used; every value of $\langle\QCD\rangle$ in
later papers is a trace over the thirty-two components. Writing
$8\pi G_{\mathrm{eff}}T_{MN}\equiv\kappa_{12}^{2}T^{\mathrm{sp}}_{MN}+\tfrac12(F_{MP}F_{N}{}^{P}-\tfrac14g_{MN}F^{2})$
and $\Lambda_{\mathrm{eff}}g_{MN}\equiv-\bigl(\Lambda_{12}+V(\QCD)\bigr)g_{MN}$
gives \eqref{eq:master}. Varying with respect to $\QCD$ gives
\eqref{eq:dilaton}. \rigR
\end{proof}

The identification of $\Lambda_{\mathrm{eff}}$ carries the dilaton's
potential, so it is constant only where $V$ is stationary; this is the
reading under which the canonical display and \eqref{eq:dilaton} are
consistent, as the next proposition shows.

\begin{proposition}[consistency]\label{prop:bianchi}
$\nabla^{M}T^{\QCD}_{MN}=\langle\Box\QCD,\nabla_{N}\QCD\rangle_{M}$. Hence the contracted
Bianchi identity, the conservation of $T_{MN}$ and \eqref{eq:master} imply
\[
\zeta\langle\Box\QCD,\nabla_{N}\QCD\rangle_{M}=\nabla_{N}V=\langle V'(\QCD),\nabla_{N}\QCD\rangle_{M}\,,
\]
which is \eqref{eq:dilaton} on the support of $\nabla\QCD$; and where
$V'=0$, $\Box\QCD=0$ there. The computation is componentwise and the trace
is linear.
\end{proposition}

\begin{proof}
Take the divergence of the stress tensor, that is, contract its first
index with a derivative. The first term, the inner product of two
gradients, gives by the product rule two pieces: the Laplacian of the
dilaton paired with its gradient, $\langle\Box\QCD,\nabla_{N}\QCD\rangle_{M}$,
and the gradient paired with the second derivative,
$\langle\nabla^{M}\QCD,\nabla_{M}\nabla_{N}\QCD\rangle_{M}$. The second term,
half the squared gradient spread over the metric, gives on differentiation
$\langle\nabla^{M}\QCD,\nabla_{N}\nabla_{M}\QCD\rangle_{M}$, the gradient
paired with the second derivative in the other order. For a scalar the two
orders of covariant differentiation agree, so those two pieces cancel and
only $\langle\Box\QCD,\nabla_{N}\QCD\rangle_{M}$ survives:
\begin{align*}
\nabla^{M}\bigl(\langle\nabla_{M}\QCD,\nabla_{N}\QCD\rangle_{M}\bigr)-\tfrac12\nabla_{N}\langle\nabla\QCD,\nabla\QCD\rangle_{M}
&=\langle\Box\QCD,\nabla_{N}\QCD\rangle_{M}\,.
\end{align*}
Now take the divergence of both sides of \eqref{eq:master}. The left side
vanishes by the contracted Bianchi identity, the geometric fact that the
Einstein tensor is automatically divergence-free. The matter term vanishes
because $T_{MN}$ is conserved (Section~\ref{sec:matter}). The vacuum term
contributes the gradient of $\Lambda_{\mathrm{eff}}$, which is the gradient
of $-V(\QCD)$, that is $-\langle V'(\QCD),\nabla_{N}\QCD\rangle_{M}$. What is
left is
\[
\zeta\langle\Box\QCD,\nabla_{N}\QCD\rangle_{M}=\langle V'(\QCD),\nabla_{N}\QCD\rangle_{M}\,,
\]
which is \eqref{eq:dilaton} wherever the dilaton is actually varying; and
where the potential is flat, $V'=0$, it says the dilaton satisfies the wave
equation $\Box\QCD=0$ there. The computation is componentwise and the trace
is linear, so nothing changes in passing from one component to the module.
\rigR
\end{proof}
%% The F1 note's "□Φ = 0 on supp ∇Φ" is the V' = 0 case; flag for the note.

\begin{proposition}[the dilaton stress tensor along the flow]\label{prop:stiff}
Let $\QCD$ depend only on the flow parameter $s$ of A4. In an orthonormal
frame $(\partial_{s},e_{1},\dots,e_{11})$ the tensor $\zeta T^{\QCD}$ is
diagonal with
\begin{equation}\label{eq:stiff}
  \zeta T^{\QCD}_{ss}=\tfrac12\zeta\langle\dot\QCD,\dot\QCD\rangle_{M}\,,\qquad
  \zeta T^{\QCD}_{ii}=-\tfrac12\zeta\langle\dot\QCD,\dot\QCD\rangle_{M}\quad(i=1,\dots,11)\,.
\end{equation}
\end{proposition}

\begin{proof}
If the dilaton depends only on the flow parameter, its gradient points
along the flow and nowhere else: $\nabla\QCD=\dot\QCD\,\partial_{s}$, where
the dot is the derivative along the flow. Its squared length is then just
$\langle\dot\QCD,\dot\QCD\rangle_{M}$. In the component along the flow the
first term of the stress tensor contributes the whole of that, and the
second term subtracts half of it, leaving a half. In any direction
transverse to the flow the first term contributes nothing, because the
gradient has no component there, and the second term subtracts half the
squared length, leaving minus a half. Off-diagonal components vanish for
the same reason (\texttt{L2\_flow\_components}). \rigR
\end{proof}

The reading. In a Lorentzian presentation, with the flow parameter
continued to a signed time, the two entries of \eqref{eq:stiff} are an
energy density and an isotropic pressure with
$\rho_{\QCD}=p_{\QCD}=\tfrac12\zeta\langle\dot\QCD,\dot\QCD\rangle_{M}$, the stiff equation of
state, and \eqref{eq:dilaton} with $V'=0$ gives $\rho_{\QCD}\propto a^{-6}$
on an isotropic background. These are the statements the programme uses.
Under Definition~\ref{def:riemannian} they are the components
\eqref{eq:stiff} read along the flow, the transverse sign carried by the
continuation; the equation is \eqref{eq:master} in either reading. \rigD

\subsection{The base restriction}\label{sec:base-restriction}

Section~\ref{sec:warp} works on the base $\Ttwo/\Ztwo$ with the fibre warped. We record
here, so that Section~\ref{sec:warp} may cite it, what the Master Equation becomes
there. Evaluate \eqref{eq:action} on the configuration of
Definition~\ref{def:scheme}: metric $g_{4}+e^{2a(\theta)}\hat g_{6}+g_{2}(\theta)$,
dilaton $\QCD=\QCD(\theta)$, the connection induced. The fibre volume
contributes $e^{6a}$, its curvature $e^{-2a}\hat R_{6}$, and the
warped-product scalar curvature contributes the gradient terms of $a$
\cite[Ch.~7]{ONeill1983}; the result is a two-dimensional action on the
base of the form
\begin{equation}\label{eq:2D}
  S_{2}\;=\;\int_{\Ttwo/\Ztwo}\!\sqrt{g_{2}}\,
  \Bigl[\frac{1}{2\kappa_{2}^{2}}R_{2}\;-\;\tfrac12\lvert F\rvert^{2}
  \;-\;\mathcal{V}(a,\QCD)\;-\;\mathcal{K}(a,\nabla a,\nabla\QCD)\Bigr]\,,
\end{equation}
with $\kappa_{2}^{-2}$ the fibre-volume-weighted gravitational
coefficient, $\mathcal{V}$ the potential assembled from $\Lambda_{12}$,
$\hat R_{6}$ and $V(\QCD)$ with the weights $e^{6a}$ and $e^{4a}$, and
$\mathcal{K}$ the gradient terms of $a$ and $\QCD$. The warp of A2 is the
geometry's response to the dilaton: on the solutions of Section~\ref{sec:warp} the two
profiles are tied by
\begin{equation}\label{eq:a-Phi}
  a\;=\;-\zeta\,\QCD\,,
\end{equation}
with $\QCD$ falling by one unit across a half-period, so that the warp
exponent across the half-period is $\zeta$ and \eqref{eq:2D} collapses to
the single-modulus Einstein--Maxwell system with potential that Section~\ref{sec:warp}
solves. That system is therefore the base restriction of
\eqref{eq:master}, not an axiom. \rigD
%% Discharges Paper 2's Σ-2 as a restriction of the Master Equation. The
%% relation eq:a-Phi is Paper 2 §7.3's definition of Φ as the modulus of the
%% warped direction, here read as the solution relation; its derivation as a
%% property of the solutions, with the boundary configuration Φ = 1, is in
%% Section 4 (WP4), which reproduces Paper 2 §4 on this footing.

\subsection{The matter sector}\label{sec:matter}

$T_{MN}$ in \eqref{eq:master} is the stress tensor of the $\Spinc$ Dirac
field of A3 on the bundle with $L=\mathcal{O}(3)$ (M1 Proposition~7.3) with
the Clifford action of the fibre algebra (M1 Definition~2.5), together
with the field strength of the induced connection. The one property of
$T_{MN}$ used in this paper is its conservation, which follows from the
diffeomorphism invariance of $S_{\mathrm{spinor}}$ on the closed manifold
$M_{12}$ (A1) by the standard argument \cite{Wald}; for the field strength
it is the Bianchi identity $\mathrm{d}F=0$ and the Maxwell equation. The
content of $T_{MN}$, its decomposition into generations and its masses,
belongs to later papers of the series and is not opened here. \rigD

\subsection{What is derived and what is read}\label{sec:summary3}

The form of \eqref{eq:master} is a theorem of the action \eqref{eq:action}
(\rigR). The action is fixed by one derivation and three readings: the
single gravitational term (minimality, \rigD), the induced connection
(\rigD), the coupling of the geometry to the sheets through $\QCD$ alone
(\rigD, with Proposition~\ref{prop:parity} as its reason), and the
Riemannian metric with time along the flow (Definition~\ref{def:riemannian}).
The coupling $\zeta$ is a constant of the theory whose value is derived in
Section~\ref{sec:warp}; the relation \eqref{eq:a-Phi} between the warp and the dilaton
is a property of the solutions established there.
